% !TeX root = ../main.tex

\section{Discussion}\label{sec:discussion}

Our study sheds light on the diverse motivations, beliefs, and mechanisms of influence that characterise crypto KOLs operating within the Web3 ecosystem. Beyond describing  individual practices, these findings carry broader implications for theories of opinion leadership as well as for the design of socio-technical systems that mediate visibility, trust, and regulation in digital finance. To situate our contributions, we structure the discussion along two perspectives: (1) theoretical implications, where we extend understandings of perceived credibility and its management under conditions of uncertainty and volatility, and (2) practical implications, where we outline design recommendations for social media platforms and supervisory authorities seeking to foster trustworthy, transparent, and community-oriented practices. We further reflect on the application of LLMs in qualitative coding and discuss limitations as well as directions for future research.

\subsection{Theoretical Implications}

One important theoretical contribution of our study is to show that the professional trajectories and day-to-day practices of crypto KOLs can be effectively interpreted through the lens of SDT. This framework has enabled us to examine how KOLs sustain long-term motivation, influence, and legitimacy within decentralised, high-stakes ecosystems characterised by shared financial exposure with followers, irreversible transactions, and thin liquidity. Thus, our findings provide evidence consistent with SDT's applicability by illustrating how the fundamental needs for autonomy, competence, and relatedness manifest in financially risky, technologically mediated, and algorithmically curated environments that differ from the domains where SDT has traditionally been applied~\cite{olafsen2025crafting,sharifi2024basic,beierlein2025creating}. Autonomy is exercised through selective sponsorships, transparent financial disclosures, explicit refusal of misaligned deals, and the rejection of promotional opportunities that could undermine credibility. Competence is demonstrated by simplifying complex knowledge, bounding claims in the face of uncertainty, producing educational content, and cultivating reputations as trusted sources of expertise. Relatedness emerges through dialogic interactions, mutual knowledge sharing, community-building practices, and maintaining social bonds across online and offline contexts.

Second, our findings reconceptualise \emph{credibility} in the Web3 influencer space. Prior research has treated finfluencer credibility as a function of static credentials, namely formal education, professional background, and financial certifications~\cite{NurhandayaniWhoDeserves2025}, reflecting broader source credibility traditions that emphasise expertise and trustworthiness as sender attributes~\cite{Metzger2013CredibilityReview,Wathen2002Credibility,Sundar2008MAIN,Flanagin2007UGC,fogg2001credibility}. Our findings challenge this view. In high-risk crypto environments characterised by widespread misinformation~\cite{Cernera2023rugpulls}, pervasive scepticism, and financial and security risks~\cite{Si2024CHI}, credibility emerges not from credentials but from \emph{self-determined, ethically enacted practices} rooted in SDT needs. We identify four dimensions of this practice-based credibility that function as community-recognised markers of trustworthiness:

\begin{enumerate}[leftmargin=*, nosep]
  \item \textbf{Self-regulation and voluntary constraint:} KOLs decline misaligned sponsorships, impose personal rules on what they will promote, and reject short-term monetisation that could erode trust (Section~\ref{sec:rq3}).
  \item \textbf{Bounded epistemic competence:} Rather than claiming omniscience, KOLs acknowledge the limits of their expertise, avoid prognostication, and emphasise that ``not even the best traders'' can reliably pick winners.
  \item \textbf{Accountability through transparency:} KOLs cultivate long-term trust through consistent disclosure of sponsorships, portfolio holdings, and conflicts of interest, accepting that transparency invites scrutiny.
  \item \textbf{Reflexive self-correction:} KOLs learn from past failures, publicly reassess their practices, and iteratively refine their norms, as illustrated by participants who abandoned paid advertisements after reputational setbacks.
\end{enumerate}

\noindent These dimensions re-conceptualize credibility as a \emph{socio-technical, self-regulated performance} rather than a static credential. Credibility is enacted through ongoing behavioural signals, reinforced by community feedback, and embedded in platform affordances (e.\,g., disclosure badges, accuracy indicators). This perspective extends source credibility theory into high-risk, decentralised contexts and offers a framework for understanding trust in environments where formal qualifications are absent, unverifiable, or irrelevant.


\subsection{Design Recommendations}

Our findings yield actionable insights for platforms, creators, and regulators seeking to foster credible communication and responsible innovation within the evolving Web3 ecosystem. While SDT illuminates why crypto KOLs pursue autonomy, competence, and relatedness, the implications here highlight how such motivations can be translated into governance, practice, and policy design.

\paragraph{Platforms: Designing for credibility and accountability.}
Platforms occupy a pivotal role in shaping socio-technical conditions for credibility. Moving beyond popularity and engagement counts, they should implement multidimensional \emph{credibility metrics} that reward transparent practices, consistent disclosures, and educational rather than sensationalist content. Practical affordances include badges for verified sponsorships or investment statements, indicators that track predictive accuracy over time, and feedback systems enabling users to flag misleading or undisclosed promotions. Equally critical is addressing the fragmentation of identity across channels: interoperable “trust profiles” could consolidate disclosure histories, accuracy scores, and community feedback into a portable reputation system. Coupled with analytical tooling—such as automated detection of undisclosed promotions, sentiment analysis, and cross-platform influence mapping—these infrastructures would make credibility more transparent and comparable, assisting users in distinguishing between hype and substantive guidance while providing structural incentives for ethical behaviour.

\paragraph{Creators and communities: Professional norms and self-regulation.}
Crypto KOLs already exhibit forms of self-regulation, including refusing misaligned promotions, labelling sponsored content, and avoiding financial fraud or similar market manipulation schemes. Formalising these practices into a community-endorsed \emph{finfluencer charter} would provide a normative benchmark for responsible engagement. Such a charter could specify disclosure formats, risk warnings, and principles for mitigating potential market impact. Beyond codification, community-governed mechanisms—peer-review panels, deliberation forums, or recognition systems led by experienced KOLs—can adjudicate contested cases, surface ethical dilemmas, and publicly acknowledge exemplary conduct. These practices would not only guide new entrants but also embed a culture of accountability that scales with the growth of Web3 communities, complementing platform infrastructures and regulatory frameworks.

\paragraph{Policy and regulation: Safeguarding users in volatile markets.}
In light of the transnational and decentralised character of crypto assets, formal oversight remains challenging, yet the potential for retail harm necessitates robust intervention. Regulators could extend existing financial promotion and consumer protection rules to encompass crypto influencers while adapting them to blockchain’s distinctive features. Measures might include standardised disclosure statements for compensated content, registration requirements for actors surpassing defined thresholds of influence, and coordinated enforcement through entities such as the International Organization of Securities Commissions (IOSCO) or regional consortia; guidance from the US Federal Trade Commission (FTC) on endorsements and testimonials, the US Securities and Exchange Commission (SEC) on crypto asset securities, the UK Financial Conduct Authority (FCA) and Advertising Standards Authority/CAP on finfluencer promotions, and the EU’s European Securities and Markets Authority (ESMA) offer relevant precedents and enforcement levers ~\cite{FTC2023EndorsementGuides, SECInvestorCrypto, FCA2023Finfluencer, ASAInfluencer2023, IOSCO2023Crypto, ESMA2023MiCAR}. Platforms can assist by identifying persistent offenders and sharing intelligence with regulators. Alongside enforcement, public agencies should invest in educational initiatives that build crypto literacy among retail investors, reducing susceptibility to misleading guidance. Taken together, harmonised regulation and proactive education would help mitigate risks while supporting responsible innovation.

Collectively, these implications extend beyond individual influencer strategies to encompass the design of accountability infrastructures, cross-platform reputation systems, community norms, and regulatory frameworks. By recognising credibility as a socio-technical outcome that demands coordinated action from platforms, creators, communities, and regulators, the Web3 ecosystem can progress toward more trustworthy and inclusive communication.

\subsection{Reflection\textit{} on LLMs}
In terms of methodology, this study is among the first to empirically demonstrate how LLMs can support qualitative coding in HCI research without displacing human oversight and interpretation. Our hybrid approach aimed to preserve interpretive rigour while reducing the cognitive and temporal load associated with labour-intensive thematic analysis. In our workflow, the LLM operated strictly upstream, suggesting additional candidate codes and SDT mappings on anonymised transcript segments for human review, rather than labelling data directly or participating in reliability calculations. Our experience shows that LLMs can assist human coders in surfacing initial themes, identifying semantic overlaps, and broadening the space of codes under consideration. Nevertheless, they remain prone to generating overly general, repetitive, or normatively loaded suggestions, justifying the need for a dual approach to coding in which humans filter, refine, and ultimately decide which labels enter the codebook. Their use also raises important questions about transparency, reproducibility, ethics, and data privacy~\cite{Schroeder2025LLMs}; we mitigated some of these concerns through strict anonymisation and human-in-the-loop validation, but model biases and versioning remain limitations. For the HCI community, these findings underscore the need for research tools and workflows that integrate LLMs responsibly, ensuring they augment rather than compromise the interpretive depth of qualitative inquiry.

\subsection{Limitations}

This study exhibits several limitations that constrain the generalisability of its findings. First, the sample comprised thirteen self-selected crypto KOLs and leaned Western, with many participants operating within European and U.S. contexts; nevertheless, the cohort also included Asian representation. Additionally, all confirmed participants were male despite active efforts to recruit female KOLs, reflecting the male-dominated composition of the crypto influencer space. Although purposive recruitment afforded heterogeneity in roles and experiences, the modest sample size, geographic skew, and lack of gender diversity restrict the extent to which the observations may be extrapolated to the wider population of crypto KOLs. Second, the reliance on semi‑structured interviews means that data are self‑reported and subject to recall and social desirability biases; participants' narratives were not cross‑validated against behavioural indicators such as on‑chain transaction records or follower engagement statistics. Third, the analysis focused exclusively on influencers' perspectives; the views of followers, regulators, and other ecosystem stakeholders were not solicited, limiting the ability to triangulate perceptions of credibility and influence. Fourth, while LLM‑assisted coding expedited early stages of thematic analysis by suggesting additional candidate labels, it may still have introduced biases inherent in large language models; despite our filtering and human verification steps, model outputs remain shaped by opaque training data and cannot substitute for human interpretive judgement. Fifth, recruitment likely exhibits survivorship and self‑selection bias: KOLs with stronger reputations, clearer disclosure norms, or more reflective practices may have been more willing to participate, potentially underrepresenting highly promotional or opaque actors. Finally, the inherent volatility of crypto markets and the rapid evolution of Web3 platforms imply that the practices and norms documented here may be transient, underscoring the need for longitudinal replication.

\subsection{Future Research Directions}

Future investigations should address these limitations and advance the field in several directions.  Comparative and cross‑cultural studies drawing on larger and more diverse samples would help to assess the universality of the present findings and reveal how regulatory and cultural contexts shape influencer behaviour.  Longitudinal mixed‑methods designs that integrate interviews with trace data—such as social media analytics, on‑chain transactions, and algorithmic amplification metrics—could illuminate how credibility develops over time and how KOLs' recommendations correlate with market dynamics.  Incorporating follower perspectives through surveys and ethnographic observation would elucidate parasocial interactions and audience responses to disclosure practices.  Methodological innovation is warranted to develop and validate computational measures of credibility, transparency, and influence that enable systematic cross‑platform assessment; such tools could draw on machine learning and natural language processing to detect undisclosed promotions and sentiment.  Examining the implications of emerging technologies—such as generative AI and decentralised identity infrastructures—for content production, reputation management, and compliance represents another important avenue.  Finally, comparative analyses across blockchain and other influencer domains could help distinguish domain‑specific behaviours from broader trends in the creator economy.

